% Rilke, "Der Panther" (1902, gemeinfrei) as a school reader page with Worterklaerungen, a biographical note and an
% Arbeitsauftrag box. German, 1 page. Features: poem with line numbers every 5 lines, glossary, task box.
\documentclass[11pt,a4paper]{article}
\usepackage[T1]{fontenc}
\usepackage[utf8]{inputenc}
\usepackage{palatino}
\usepackage[scaled=0.9]{helvet}
\usepackage[ngerman]{babel}
\usepackage[margin=2.0cm,top=2.1cm,bottom=2.0cm]{geometry}
\usepackage{xcolor}
\usepackage{tcolorbox}
\usepackage{enumitem}
\usepackage{fancyhdr}
\usepackage{microtype}

\definecolor{lesebuch}{HTML}{1F5A7A}
\definecolor{lesebuchhell}{HTML}{E4EFF5}
\setlength{\parindent}{0pt}
\setlength{\parskip}{0.5em}

\pagestyle{fancy}
\fancyhf{}
\renewcommand{\headrulewidth}{0.4pt}
\fancyhead[L]{\sffamily\small\textcolor{lesebuch}{\textbf{Deutsch Lesebuch 9/10}} \textperiodcentered{} Lyrik}
\fancyhead[R]{\sffamily\small Kapitel 4: Dinggedichte}
\fancyfoot[L]{\sffamily\footnotesize Text gemeinfrei (Rilke, gest.\ 1926). Aufgaben und Erläuterungen: eigener Text für Testzwecke.}
\fancyfoot[R]{\sffamily\footnotesize 112}

\begin{document}

% ---- Kapitelkopf ------------------------------------------------------------------------
{\sffamily\small\color{lesebuch}\textbf{TIERE IN DER LYRIK}}\\[-0.2em]
{\color{lesebuch}\rule{\linewidth}{0.8pt}}

\vspace{0.4em}
\begin{minipage}[t]{0.63\linewidth}\vspace{0pt}
% ---- Gedicht ----------------------------------------------------------------------------
{\Huge\bfseries Der Panther}\\[0.2em]
{\large\itshape Im Jardin des Plantes, Paris}\\[0.5em]
{\sffamily\small Rainer Maria Rilke (1875--1926)}

\vspace{1.4em}
\large
\renewcommand{\arraystretch}{1.28}
\begin{tabular}{@{}r@{\hspace{1.6em}}l@{}}
   & Sein Blick ist vom Vorübergehn der Stäbe \\
   & so müd geworden, daß er nichts mehr hält. \\
   & Ihm ist, als ob es tausend Stäbe gäbe \\
   & und hinter tausend Stäben keine Welt. \\[0.9em]
  \textcolor{lesebuch}{\sffamily\footnotesize 5} & Der weiche Gang geschmeidig starker Schritte, \\
   & der sich im allerkleinsten Kreise dreht, \\
   & ist wie ein Tanz von Kraft um eine Mitte, \\
   & in der betäubt ein großer Wille steht. \\[0.9em]
   & Nur manchmal schiebt der Vorhang der Pupille \\
  \textcolor{lesebuch}{\sffamily\footnotesize 10} & sich lautlos auf --. Dann geht ein Bild hinein, \\
   & geht durch der Glieder angespannte Stille -- \\
   & und hört im Herzen auf zu sein. \\
\end{tabular}
\normalsize

\vspace{1.0em}
{\sffamily\footnotesize Aus: Rainer Maria Rilke, \emph{Neue Gedichte}, 1907. Das Gedicht entstand 1902 in Paris.}
\end{minipage}\hfill
\begin{minipage}[t]{0.33\linewidth}\vspace{0pt}
% ---- Randspalte -------------------------------------------------------------------------
\begin{tcolorbox}[colback=lesebuchhell, colframe=lesebuch, boxrule=0.5pt, arc=2pt, left=6pt, right=6pt,
  top=4pt, bottom=4pt, title={\sffamily\bfseries Worterklärungen}, fonttitle=\sffamily\small, coltitle=white,
  colbacktitle=lesebuch]
\sffamily\footnotesize
\textbf{Jardin des Plantes} (frz.): botanischer Garten mit Tierpark in Paris\\[0.2em]
\textbf{geschmeidig:} biegsam und elegant in der Bewegung\\[0.2em]
\textbf{betäubt:} gelähmt, ohne Gefühl\\[0.2em]
\textbf{Pupille:} das schwarze Sehloch in der Mitte des Auges\\[0.2em]
\textbf{angespannt:} gespannt, unter Kraft gehalten
\end{tcolorbox}

\vspace{0.3em}
\begin{tcolorbox}[colback=white, colframe=lesebuch, boxrule=0.5pt, arc=2pt, left=6pt, right=6pt,
  top=4pt, bottom=4pt, title={\sffamily\bfseries Zum Autor}, fonttitle=\sffamily\small, coltitle=white,
  colbacktitle=lesebuch]
\sffamily\footnotesize
Rainer Maria Rilke wurde 1875 in Prag geboren und starb 1926 in der Schweiz. 1902 reiste er nach Paris, um
über den Bildhauer Auguste Rodin zu schreiben. Rodin lehrte ihn, die Dinge lange und genau anzusehen. So
entstand eine neue Art von Gedicht, das \emph{Dinggedicht}: Es beschreibt einen Gegenstand oder ein Tier
so genau, dass es zum Sinnbild wird. \emph{Der Panther} gilt als das bekannteste Beispiel.
\end{tcolorbox}
\end{minipage}

% ---- Arbeitsauftrag ----------------------------------------------------------------------
\vspace{1.2em}
\begin{tcolorbox}[colback=white, colframe=lesebuch, boxrule=1.2pt, arc=3pt, left=10pt, right=10pt, top=6pt, bottom=6pt,
  title={\sffamily\bfseries Arbeitsauftrag \hfill\normalfont\small Bearbeitungszeit: 25 Minuten}, fonttitle=\sffamily\large,
  coltitle=white, colbacktitle=lesebuch]
\begin{enumerate}[leftmargin=1.6em, itemsep=0.35em, label=\textbf{\arabic*.}]
  \item Lies das Gedicht zweimal laut. Halte deinen ersten Eindruck in einem Satz fest.
  \item \textbf{Analysiere das Reimschema und das Metrum} des Gedichts. Bestimme auch die Kadenzen
        (männlich oder weiblich) der Verse.
  \item Untersuche die \textbf{sprachlichen Bilder} (zum Beispiel Vergleiche, Metaphern und
        Übertreibungen) und erkläre ihre Wirkung.
  \item Beschreibe, wie sich der Zustand des Panthers von Strophe zu Strophe verändert. Belege deine Aussagen am
        Text und gib die Zeilen an.
\end{enumerate}
\end{tcolorbox}

\end{document}
